\section{Background}

\label{sec:background}
%%==========================================================================

\subsection{Lux Appchain Architecture}

A Lux appchain consists of:
\begin{itemize}[leftmargin=1.5em]
  \item A subset of Primary Network validators who additionally validate the appchain;
  \item A custom virtual machine (VM) specifying the state transition function;
  \item An independent consensus instance (Beacon or custom);
  \item Optional: native gas token, custom fee structure, validator admission policy.
\end{itemize}

Appchains are registered on the P-Chain via a \texttt{CreateAppchain} transaction, which
requires a one-time fee of 1 LUX (as of this writing). Validators opt into appchain
validation via \texttt{AddAppchainValidator} transactions. Critically, appchain validators
are not required to use LUX as their gas token; they may denominate fees in any ERC-20
token or in the appchain's native currency.

\subsection{Economic Primitives}

Let $\mathcal{N}$ denote the set of Primary Network validators with staking weights
$\{w_v\}_{v \in \mathcal{N}}$ and total stake $W = \sum_v w_v$. The annual Primary Network
staking yield for validator $v$ is:
\[
  y_v = \frac{\text{AnnualPrimaryReward}(v)}{w_v \cdot P_{\text{LUX}}},
\]
where $P_{\text{LUX}}$ is the LUX price in USD. This yield represents the \emph{opportunity
cost} that appchain incentives must exceed to attract validators.

\subsection{Related Work}

Cosmos Zone economics~\cite{kwon2016cosmos} are the closest analogue: application-specific
chains with sovereign validator sets secured by ATOM-denominated staking. Key differences
from Lux appchains include: Cosmos chains maintain independent validator sets (no sharing),
IBC is the interoperability layer (versus Lux's native cross-chain messaging), and ATOM
inflation subsidizes security rather than fee revenue.

Ethereum Layer 2 economics~\cite{konstantopoulos2021optimism} address a different
problem---reducing mainnet congestion through off-chain execution with on-chain validity
guarantees---but share the challenge of bootstrapping sequencer economics without
centralization.

Polkadot parachain auctions~\cite{wood2016polkadot} provide an interesting comparison:
parachains bid for relay-chain security through DOT bond auctions, creating a market-driven
mechanism for security allocation. LRP appchains instead pay validators directly, which is
more flexible but requires more explicit incentive design.

%%==========================================================================
